\section*{Βήμα 14}
\subsection*{1}
Η υλοποίηση του βήματος αυτού βασίζεται στη δομή του βοηθητικού κώδικα για το εργαστήριο και συγκεκριμένα στο αρχείο \textit{lstm.py}. Έχουν συμπληρωθεί τα τμήματα σημασμένα ως TODO και επιπλέον έχουν τροποποιηθεί λειτουργικότητες και προστεθεί επιπλεόν συναρτήσεις ώστε να καλυφθούν τα υποερωτήματα του βήματος.

\subsection*{2}
Για την κατασκευή του LSTM δικτύου αξιοποιήσαμε τη βιβλιοθήκη PyTorch. Ο κώδικας φαίνεται παρακάτω:

\lstinputlisting[language=Python, firstline=29, lastline=49]{../../lab.py}

\noindent Κατά την αρχικοποίηση του δικτύου ελέγχεται κατά πόσο η αρχιτεκτονική του είναι αμφίδρομη ή όχι. Εφόσον είναι, αφενός τίθεται σε True η αντίστοιχη παράμετρος και αφετέρου φροντίζουμε ώστε να διπλασιαστεί το μέγεθος του hidden state του δικτύου. \\

\noindent Στη συνέχεια, επισημαίνουμε τα είδη των layers που θα χρησιμοποιηθούν στην υλοποίηση του δικτύου. Το κυρίως layer είναι το nn.LSTM, το οποίο ρυθμίζεται με τις παραμέτρους του με τρόπο τέτοιο ώστε να καλύψει τις απαιτήσεις της εκφώνησης στα πιθανά ερωτήματα που θα προκύψουν.
\noindent Το dropout layer δεν χρησιμοποιείται πάντα. Αυτό το καθορίζουμε με την παράμετρο dropout του nn.Dropout layer. Εάν αυτή είναι 0 προφανώς αυτή η τεχνική regularization δεν χρησιμοποιείται. Επιπλέον, αξίζει να αναφέρουμε ότι έχουμε φροντίσει το Dropout να μην εφαρμοστεί στο τελευταίο layer.
\noindent Το nn.Linear layer χρησιμοποιείται ώστε να πάρει τα output του LSTM (για αυτό έχει και μέγεθος όσο του hidden state) και να τα συνδυάσει προκειμένου να πάρει την τελική απόφαση για το classification του ψηφίου. \\

\noindent Η συνάρτηση forward είναι αυτή που καθορίζει το πώς γίνεται το inference του μοντέλου. Δηλαδή, κατ' επέκταση, τη διάταξη των δομικών layer που ορίσαμε και εξηγήσαμε παραπάνω.
\noindent Ο πρώτος διαχωρισμός που κάνουμε έγκειται στο αν χρησιμοποιούμε την τεχνική packing στο input (χρειάζεται για το ερώτημα 8 - Bonus). Με την τεχνική του padding πετυχαίνουμε 2 πράγματα. Αρχικά, σε dataset με δεδομένα ακολουθίες διαφορετικών μηκών, είναι σκόπινο να ενώσουμε τις μιρκότερες από αυτές ώστε να τις επεξεργαζόμαστε ταυτόχρονα και να εξοικονομούμε χρόνο κατά την εκπαίδευση. Επιπρόσθετα, αποφεύγουμε το zero-padding και την επεξεργασία αυτής της παραμορφωμένης εκδοχής των ακολουθιών εισόδου. Αυτή η τεχνική είναι ιδιαίτερα χρήσιμη σε δίκτυα στα οποία το μέγεθος εισόδου είναι αυστηρά καθορισμένο, όπως οι Transformers. 
\noindent Ο τρόπος που επιλέγουμε να εκτελέσουμε το packing είναι αρχικά να ταξινομήσουμε τις ακολουθίες με τα φωνήματα κατά φθίνον μήκος έτσι ώστε να τις συνδυάσουμε αποτελεσματικότερα. Αυτός ο συνδυασμός πραγματοποιοείται από την συνάρτηση \textit{packed\_padded\_sequence()} της βιβλιοθήκης της Pytorch (βλέπουμε ότι το enforced\_sorted είναι ορισμένο σε True). \\

\noindent Τέλος, όπως βλέπει κανείς εδώ:

\lstinputlisting[language=Python, firstline=492, lastline=496]{../../lab.py}

\lstinputlisting[language=Python, firstline=499, lastline=513]{../../lab.py}

\noindent Έτσι εξασφαλίζουμε ότι το dropout δε θα μπει στο τελευταίο layer. 

\subsection*{3}
Σε αυτό το βήμα καλούμαστε να εκπαιδεύσουμε το μοντέλο. Έτσι, αρχικά, θα εξηγήσουμε τη διαδικασία της εκπαίδευσής του. \\

\lstinputlisting[language=Python, firstline=554, lastline=581]{../../lab.py}

\noindent Στο trainning loop παίρνουμε τα δεδομένα του dataloader batch by batch, υπολογίζουμε το loss και κάνουμε το ανάλογο βήμα στον optimizer. \\

\noindent Ο optimizer που χρησιμοποιούμε είναι ο AdamW. Αξιοποιεί όρους ορμής και έχει μοντέλο για να μεταβάλει τους όρους του learning rate και της ταχύτητας κατάλληλα ώστε να επιτύχει την καλύτερη σύγλιση. Προφανώς εκτυπώνουμε το loss σε κάθε εποχή όπως μας ζητείται στην εκφώνηση. 

\subsection*{4}

Στην συνάρτηση train:

\lstinputlisting[language=Python, firstline=728, lastline=774]{../../lab.py}

\noindent Παίρνουμε τις μετρικές που μας ενδιαφέρουν, δηλαδή το trainning και validation loss καθώς και το validation accuracy. Η master train είναι αυτή που εκπαιδεύει όλα τα μοντέλα που μας ενδιαφέρουν.

\subsection*{5}
Καλούμαστε να προσθέσουμε L2 Regularization στο μοντέλο μας, καθώς και dropout. \\

\noindent Το L2 regularization δεν είναι κάτι το οποίο επηρεάζει δομικά το μοντέλο μας άμεσα. Είναι ένας όρος ποινικοποίησης του μέτρου των βαρών που προστίθεται στο loss function. Έτσι, μειώνουμε την πολυπλοκότητα του μοντέλου και έχουμε ως αποτέλεσμα την αποφυγή του overfitting. Αυτό είναι σημαντικό και αναμένεται να παίξει ρόλο, ειδικά σε ένα απλό dataset όπως είναι το δικό μας. 


\begin{itemize}
	\item epochs: 750
	\item learning rate: $2*10^-5$
	\item dropout: 0.4
\end{itemize}
